\subsection{Design principles}
\label{sec:principles}

Well-formedness enforces several constraints, summarised here with the programs each accepts and rejects;
the rules follow in \secref{well-formedness-statements} and \secref{well-formedness-modules}.

\begin{DeletedBlock}
\subsubsection{Assignment is definition}

Python has no variable declarations: a variable is introduced by assigning to it, and a local variable is
in scope throughout the function body, wherever the assignment occurs. \PurePy keeps the syntax but gives
each reference a unique lexical referent through \emph{definite assignment}: a variable may be referenced
only where an assignment covers every path to the reference. The \ruleName{var} rule requires the variable
to be definitely assigned, so that $\Gamma(x)$ is a type; the \ruleName{assign} rules make it so, with
$x = e$ assigning $\set{\bind{x}{\sigma}}$ for the type $\sigma$ of $e$; and the \ruleName{if} rules
combine the branches with $\merge$, so that a conditional definitely assigns a variable only if every
branch that does not definitely return assigns it, otherwise leaving it at $\False$
(\secref{well-formedness-statements}). Branch-local
variables are therefore permitted, but using one after the conditional is not:

\begin{python}
def foo():
    x = 6
    if b:
        x = 3
        y = x + 1
    else:
        y = 0
    return x  # error: partially defined
\end{python}

\noindent Python returns \pyinline{3} when \pyinline{b} is true, so no pure reading of the program agrees
with Python. The program assigns \pyinline{x} on every run, so this is stricter than avoiding Python's
\pyinline{NameError}: the \pyinline{x} in \pyinline{return x} would refer to the branch-local variable on
one path and to the variable assigned by \pyinline{x = 6} on the other.

Python's function-level scoping is kept as well. The \ruleName{def} rule introduces every variable assigned by
the body, $\assignsS(b)$, at $\False$ before checking the body, so a reference before the
assignment is rejected statically, where Python fails at run time with \pyinline{UnboundLocalError}:

\begin{python}
y = 7
def f():
    x = y  # error: y is local, because of the assignment below
    y = 8
\end{python}

\noindent Here $\assignsS(s)$ is the set of variables assigned anywhere in a statement $s$, not descending
into nested function definitions, which have their own scope.

\subsubsection{No reassignment of captured variables}

In Python a closure holds a mutable binding, so a variable captured by a lambda or nested function
definition may be reassigned afterwards, changing the value seen by the function (\secref{no-mutable-state}).
\PurePy excludes this: a variable captured by a statement may not be assigned later in the same sequence, nor
by that statement itself. The \ruleName{seq} rule imposes the first as a side condition on a sequence
$t\;t'$: no variable in $\captures(t)$ may be in $\assignsS(t')$. In the example of \secref{no-mutable-state}, \pyinline{x} is captured by the
definition of \pyinline{g} and assigned by \pyinline{x = 6}, so the program is rejected. The
\ruleName{assign} rule imposes the second, $x \notin \capturesE(e)$ for an assignment $x = e$, rejecting
for instance \pyinline{x = lambda: x}, where under Python's late binding the lambda would return itself. A
third case arises in comprehensions: in Python a comprehension mutates its generator variables on each
iteration, so a lambda in the body or in a later qualifier may not capture one, a side condition of
\ruleName{list-comp}, \ruleName{dict-comp} and \ruleName{qual-generator}. For instance
\pyinline{[lambda x: x**i for i in [1, 2, 3]]} is rejected, since under Python each lambda in the list would compute
\pyinline{x**3}. A dataclass declaration counts as an assignment of its class name here, although re-declaring a class name is
excluded outright (\secref{no-first-class}).

Here $\captures(s)$ is the set of variables of the enclosing scope captured by the lambdas and nested
function definitions within a statement $s$, with $\capturesE(e)$ computing the same for an expression $e$: the
free variables of each such function's body, less its parameters and any variables assigned by the body, which
are local to it. For a mutual region of function definitions (\figref{syntax-stmt}) the names of the functions
themselves are also excluded, since the \ruleName{def} rule binds them in every body, consistent with
Python's late binding semantics (\secref{mutual-recursion}).

\end{DeletedBlock}
\begin{AddedBlock}
\subsubsection{Declaration and definite assignment}
\label{sec:local-assignment}

Python has no variable declarations: a variable is introduced by assigning to it, and a local variable is
in scope throughout the function body, wherever the assignment occurs. \PurePy requires a declaration, an annotation of the name such as \pyinline{x: int}, which Python has allowed since version 3.6, and separates
declaring a variable from assigning it through \emph{definite assignment}: a variable may be referenced only
where an assignment covers every path to the reference, and assigned only where no path has assigned it
already. The \ruleName{declare} rule binds the variable at $\DU{\tau}$, declared at $\tau$ and
definitely unassigned; the \ruleName{assign} rule requires that entry, checks the expression against $\tau$ and
makes the variable definitely assigned, so that $\Gamma(x)$ is a type, as the \ruleName{var} rule requires; an
annotated assignment does both at once; and the \ruleName{if} rules combine the branches with $\merge$, so
that a conditional definitely assigns a variable only if every branch that does not definitely return
assigns it, otherwise leaving it possibly unassigned, $\PU{\tau}$ (\secref{well-formedness-statements}). A variable is thus
declared once in its scope and assigned at most once on any path, following Java's rule for blank finals~\cite[\S16]{gosling2025}. Branch-local variables are therefore permitted, but using one after the
conditional is not, nor is completing a partial assignment:

\begin{python}
def foo(b: bool) -> int:
    x: int
    if b:
        x = 3
        y: int = x + 1
        print(y)
    x = 0     # error: x assigned on the path through the branch
    return x
\end{python}

\noindent Python assigns \pyinline{x} twice when \pyinline{b} is true, so no pure reading of the program
agrees with Python: under shadowing, \pyinline{return x} would refer to a different variable from the one
assigned in the branch. Following mypy, a name may be declared only once in a scope, so a second annotation
of the same name is rejected even in another branch of a conditional, as is a declaration of a parameter, of
an imported name or of a name bound by a pattern.

Python's function-level scoping is kept as well. The \ruleName{def} rule introduces every variable declared
or assigned by the body, $\assignsS(s)$, at $\Unbound$ before checking the body, so a reference before the
declaration is rejected statically, where Python fails at run time with \pyinline{UnboundLocalError}:

\begin{python}
y: int = 7
def f() -> int:
    x: int = y  # error: y is local, because of the declaration below
    y: int = 8
    return x
\end{python}

\noindent Here $\assignsS(s)$ is the set of variables declared or assigned anywhere in a statement $s$, not
descending into nested function definitions, which have their own scope.

\subsubsection{No capture of comprehension variables}
\label{sec:captured-variables}

In Python a closure holds a mutable binding, so a variable captured by a lambda or nested function
definition may be reassigned afterwards, changing the value seen by the function (\secref{no-mutable-state}).
This cannot arise in \PurePy: a closure can capture only a definitely assigned variable, and no later
statement may assign one. One case remains, in comprehensions: in Python a comprehension mutates its generator
variables on each iteration, so a lambda in the body or in a later qualifier may not capture one, a side
condition of \ruleName{list-comp}, \ruleName{dict-comp} and \ruleName{qual-generator}. For instance
\pyinline{[lambda x: x**i for i in [1, 2, 3]]} is rejected, since under Python each lambda in the list would
compute \pyinline{x**3}.

Here $\captures(\seq{g})$ is the set of variables of the enclosing scope captured by the lambdas within the
qualifiers $\seq{g}$ of a comprehension, with $\captures(e)$ computing the same for an expression $e$: the
free variables of each lambda's body, less its parameters.

\end{AddedBlock}
\subsubsection{Mutual regions}
\label{sec:mutual-recursion}

Mutual recursion in Python relies on late binding: the free variables of a function are resolved when
they are evaluated, so two functions may refer to each other whatever order they are defined in, provided
both exist by the time the reference is evaluated. \PurePy makes the
dependency static: a \emph{mutual region}, a maximal run of consecutive \pyinline{def} statements in a
sequence (\figref{syntax-stmt}), is a single simultaneous binding. The \ruleName{def} rule checks each body in a
context in which every name of the mutual region is definitely assigned, so any function in the region may refer
to any other, and requires the names to be distinct. Self-recursion is the special case of a singleton
mutual region. Any other statement between two definitions separates them into different mutual regions:

\begin{DeletedBlock}
\begin{python}
def even(n):
    if n == 0:
        return True
    return odd(n - 1)  # error: odd is not yet bound

x = even(0)  # intervening statement separates the mutual regions

def odd(n):
    if n == 0:
        return False
    return even(n - 1)
\end{python}
\end{DeletedBlock}
\begin{AddedBlock}
\begin{python}
def even(n: int) -> bool:
    if n == 0:
        return True
    return odd(n - 1)  # error: odd is not yet bound

x: bool = even(0)  # intervening statement separates the mutual regions

def odd(n: int) -> bool:
    if n == 0:
        return False
    return even(n - 1)
\end{python}
\end{AddedBlock}

\noindent This is rejected statically: \pyinline{even} forms a mutual region by itself, in which \pyinline{odd}
is not yet bound, so the reference fails the \ruleName{var} rule. Python accepts the program, since the
call \pyinline{even(0)} never reaches \pyinline{odd}. A mutual region with two definitions of the same name is
rejected by the distinct-names condition:

\begin{DeletedBlock}
\begin{python}
def g():
    return 0

def f():
    return g()
def g():     # error: g defined twice in the same mutual region
    return 1
\end{python}
\end{DeletedBlock}
\begin{AddedBlock}
\begin{python}
def g() -> int:
    return 0

def f() -> int:
    return g()
def g() -> int:  # error: g defined twice in the same mutual region
    return 1
\end{python}
\end{AddedBlock}

\noindent Python accepts this too, treating the second \pyinline{def g} as a reassignment of \pyinline{g},
which late binding then makes the target of the call in \pyinline{f}. Finally, the branches of a conditional may contain mutual regions, with any use of those
functions after the conditional subject to the usual definite assignment rules (\secref{local-assignment}).

\subsubsection{Pattern binding and reachability}
\label{sec:pattern-binding}

A \pyinline{case} pattern binds variables, and Python scopes them as it does assignments: they remain in
scope after the \pyinline{match} statement. \PurePy follows: each case is treated as a branch, and its pattern as an
assignment of the variables it binds\Added{, which must not be declared already (\ruleName{pat-var})}. The set of variables a pattern $p$ introduces is $\binds(p)$; the
\ruleName{match} rule checks each body in a context in which they are definitely assigned, and combines
the cases with $\merge$ as the \ruleName{if} rules do, so a variable is definitely assigned after the match
only if every case binds it or returns. A match whose cases leave some value of the scrutinee type
unmatched may fall through; in this case the \ruleName{match-partial} rule merges in the empty assignment
set, as for a conditional without an \pyinline{else}:

\begin{DeletedBlock}
\begin{python}
v = (1, 2)
match v:
    case (x, y):
        print(x, y)
print(x)  # error: x not definitely assigned
\end{python}
\end{DeletedBlock}
\begin{AddedBlock}
\begin{python}
v: tuple[int, int] = (1, 2)
match v:
    case (x, y):
        print(x, y)
print(x)  # error: x not definitely assigned
\end{python}
\end{AddedBlock}

\noindent Python prints \pyinline{1}, the binding having escaped the match. Within a pattern the
variables must be distinct, as in Python. Every case must also be reachable: the matching judgement rejects a case
that can match only values already matched by an earlier case (\secref{well-formedness-patterns}), which
Python leaves silently unreachable:

\begin{DeletedBlock}
\begin{python}
v = (1, 2)
match v:
    case (a, b):
        print("seq", a, b)
    case (1, 2):  # error: the previous case matches every pair
        print("lit")
\end{python}
\end{DeletedBlock}
\begin{AddedBlock}
\begin{python}
v: tuple[int, int] = (1, 2)
match v:
    case (a, b):
        print("seq", a, b)
    case (1, 2):  # error: the previous case matches every pair
        print("lit")
\end{python}
\end{AddedBlock}

\subsubsection{No unreachable code}
\label{sec:unreachable}

Python accepts statements that can never run, such as code after a \pyinline{return}. \PurePy rejects
them, the statement-level counterpart of case reachability (\secref{pattern-binding}). A statement
sequence returns only if its last statement returns and every earlier one assigns: the \ruleName{seq}
rule requires the first statement of a sequence to have static outcome $\Assigns{\Delta}$
(\secref{well-formedness-statements}), so a statement following a \pyinline{return} has no derivation:

\begin{DeletedBlock}
\begin{python}
def foo(x):
    return x + 1
    return x  # error: unreachable
\end{python}
\end{DeletedBlock}
\begin{AddedBlock}
\begin{python}
def foo(x: int) -> int:
    return x + 1
    return x  # error: unreachable
\end{python}
\end{AddedBlock}

\subsubsection{Static module loading}
\label{sec:load-order}

In Python a module's namespace is populated as loading proceeds, by the assignments in its body and, for
a module with submodules, by the import system binding each submodule as it loads, so the referent of a qualified name can
depend on which modules have loaded so far, anywhere in the program (\secref{no-load-dependent}). \PurePy gives each name a meaning fixed by the
text and imports of the module alone. Imports come first in a module, so its dependencies are known before any
statement runs. Loading is checked statically: the \ruleName{module} rule checks a body and computes its
signature, and the \ruleName{import} and \ruleName{from-import} rules statically load each target and its
containing modules first, so an import cycle has no derivation. A module may not bind a name that is also the name of one of its
submodules, a side condition of \ruleName{module}, since Python resolves the clash by load order. A
qualified reference must resolve through the imports of the referring module to a definitely assigned
member (\ruleName{attr-module}). A submodule is not an attribute of its parent until it is loaded, so
\pyinline{import pkg} followed by \pyinline{pkg.sub.y} is rejected unless the module also imports
\pyinline{pkg.sub}, whereas Python permits the reference whenever some earlier import, anywhere in the
program, happens to have loaded \pyinline{pkg.sub}:

\begin{minipage}[t]{0.3\textwidth}
\begin{python}
# a
import pkg.sub
\end{python}
\end{minipage}
\hfill
\begin{minipage}[t]{0.6\textwidth}
\begin{python}
# __main__
import pkg
import a
print(pkg.sub.y)  # error: pkg.sub is not imported here
\end{python}
\end{minipage}

\noindent Python accepts this because importing \pyinline{a} loads \pyinline{pkg.sub}; remove that import,
which the main module does not otherwise use, and Python raises \pyinline{AttributeError}. It raises the
same for a reference to a module member never assigned by the body, such as \pyinline{lib.x} where
\pyinline{lib} assigns \pyinline{x} only on a branch that is not taken. \PurePy rejects all three programs
statically (\secref{well-formedness-modules}).

